\section{Testing a column with many symbols}\label{sec:design}

In the planar proof an implication is tested for every possible pair
of equal symbols. A column with \(r\) entries has \(\binom r2\) such
pairs. To obtain the bound of Theorem~\ref{thm:main} with its stated
dependence on \(r\), we construct a single test for the whole column.
It returns a fixed nonzero tensor on agreement and vanishes on every
\emph{violation}, meaning that some symbol repeats but the column is
nonconstant. Its value on pairwise distinct symbols may be arbitrary.

Take \(r\ge4\). Write \(H=\log(2r)\) and \(L=\log(2H)\), with all
logarithms in base two. Two of the symbols are held on groups of sites
\(A,B\), called the anchors. The other symbols are held at individual
satellite sites. The recursive test will sometimes use only a subset
of the satellites, called the active satellites. The two anchors have
joint width one and individual widths summing to \(1-s\). Since the
satellites are points, this leaves a surplus \(s>0\) in the word-count
comparison of Section~\ref{sec:words}. We seek a bound on the cost of
one column that is independent of the length of its symbols.

\paragraph{Shortening a symbol.}
Split each symbol into a head \(a\) and a tail \(v\) of \(b'\) bits.
For every possible head, sum over a mask \(m_a\) of \(b'\) bits shared
by all sites. Apply the shorter test to the column of shifted tails
\(v+m_a\), with bitwise addition modulo two. Retain the masks in the
output. On agreement every choice of masks survives, so the test
returns their full diagonals. On a violation the repeated symbol stays
repeated. The shorter test can therefore survive only if all the
shifted tails take a single value \(z\). We have to subtract the masks
producing this \emph{false tie}.

For a given \(z\), every entry \((a,v)\) forces \(m_a=v+z\). If two
tails at one head disagree, there is no solution. An empty head
leaves its mask free. Figure~\ref{fig:example} gives a five-entry
column in which only the two satellites \(P,R\) repeat a letter.
At \(z=01\), the two occurrences of head \(10\) both require
\(m_{10}=10\).

\begin{figure}[!htbp]
\centering
\begin{tikzpicture}[x=1cm,y=1cm,font=\small,
 lab/.style={anchor=east,inner sep=0pt},bit/.style={font=\small\ttfamily},
 line/.style={draw=ink!45},child/.style={draw=rust!75,fill=rust!5,
 rounded corners=2pt,inner xsep=6pt,inner ysep=4pt}]
 \fill[sea!7,rounded corners=2pt] (2.26,-.32) rectangle (4.82,2.28);
 \draw[sea!75,rounded corners=2pt] (2.26,-.32) rectangle (4.82,2.28);
 \foreach \y/\name in {2/Site,1.49/Head,.98/Tail,.47/Mask for \(z=01\),-.04/Shifted tail}
   \node[lab] at (-.7,\y){\name};
 \foreach \x/\site/\head/\tail/\mask in
   {0/A/00/10/11,1.35/B/01/01/00,2.87/P/10/11/10,4.21/R/10/11/10,5.63/T/11/10/11}{
    \node at (\x,2){\(\site\)};
    \node[bit] at (\x,1.49){\head};
    \node[bit] at (\x,.98){\tail};
    \node[bit,text=sea] at (\x,.47){\mask};
    \node[bit] at (\x,-.04){01};
 }
 \draw[line] (-.43,1.76)--(6.06,1.76);
 \draw[line] (-.43,.22)--(6.06,.22);
 \draw[-{Latex[length=2mm]},thick,sea] (3.54,-.32)--(3.54,-.7);
 \node[child] at (3.54,-1.02) {\(\bigl(u,t,\mathtt{10},\mathtt{10}\bigr)\)};
 \node[lab] at (2.14,-1.02){Shorter column};
\end{tikzpicture}
\caption{One false tie and its correction. The shaded satellites carry
the same original symbol. Their required mask becomes their input in
the shorter column, alongside the two auxiliary anchor inputs \(u,t\).
The repeat forces this column to pass only at \(u=t=\mathtt{10}\).
The sum over \(u,t\) thus recovers the required mask.}
\label{fig:example}
\end{figure}

\paragraph{Recovering the cost of a false tie.}
We could supply every mask, sum over \(z\), and check the constraints
\(m_a=v+z\) locally. Relative to the mask diagonals returned on
agreement, this would leave \(b'\) bits to pay for \(z\). First
consider a violation in which neither anchor symbol repeats. At a
false tie each anchor must be alone in its head, since every entry in
a given head must have the same tail. Each anchor can therefore
compute its mask locally from \(z\) and broadcast it. A broadcast
from \(A\) to all sites costs \((1-w(A))b'\), saving \(w(A)b'\) against
the full mask diagonal. The two anchors together pay for
\((1-s)b'\) of the cost of \(z\). This leaves \(sb'\).

A head containing a genuine satellite repeat can cover the remainder.
Give two auxiliary anchor groups independent dummy labels \(u,t\);
their widths sum to \(1-s'\), with \(s'\ge s\). For each active
satellite in the chosen head, give the shorter test its required mask
\(v+z\). If a whole symbol repeats in that head, the shorter column
has a repeat for every pair \(u,t\). If its required masks differ, it
always vanishes. If the masks agree, exactly one pair \(u,t\) makes
the shorter column constant. The auxiliary anchors can then broadcast
their common label from their joint width of one. The cost of the two
dummy labels saves \(s'b'\) against the full mask diagonal, before
the cost of the shorter test itself.

Inclusion--exclusion removes the need to identify the repeated head. Let
\(I_i\) be the exact mask operation for head \(i\), and let \(T_i\) be
the operation using the shorter test, with the auxiliary tensor
factors made common to the terms. With two possible non-anchor heads,
use
\[
          I_0\otimes I_1-(I_0-T_0)\otimes(I_1-T_1).
\]
At a repeated symbol in head \(i\) we have \(T_i=I_i\), so the
display equals the exact mask correction. When the product is
expanded, every remaining term uses at least one shorter test.
Section~\ref{sec:masks} carries out this cancellation for any number
of heads, keeping each mask tuple in its own output slot.

\paragraph{Controlling the size of the shorter test.}
A preliminary gate, the \emph{guard}, handles violations with a
repeated anchor symbol. It also kills a violation when one symbol
occurs at more than \((1-p_0)n\) of the \(n\) active satellites, for
a fixed \(p_0>0\). On a surviving violation a head larger than this
threshold cannot have a common tail. Its exact correction is then
zero, and we can omit the shorter test at that head. Thus a recursive
call that changes the active satellites always reduces their number
by a fixed fraction.

The difficulty in implementing the guard is to compare a widely
repeated symbol with the anchors without sending every satellite
symbol to them. On the guard's large-majority promise, a balanced split of the active
satellites into two geometrically separated classes leaves the same
symbol in a strict majority in each class. For any projection of the letters onto \(u\)
bits, Section~\ref{sec:vote} sends the two projected majorities to
two designated destinations, charging one wire per class. The
placement in Section~\ref{sec:geometry} makes the width of the two
destinations exceed the combined wire rates by more than \(s\).
Together with the anchor widths \(1-s\), this gives the strict width
inequality in Lemma~\ref{lem:oracle}. That lemma recovers the equality
test for full symbols from tests of their projections. Once both
anchors agree, the test can check every active satellite locally.
Section~\ref{sec:guard} combines these operations.

Shortening the letters alone preserves the split. We pay to determine
a new split only when a correction changes the active satellites.
For the binary layout used in the proof, there are \(O(\log r)\) such
changes along any branch. In
Section~\ref{sec:recurrence} this yields first a linear cost with
slope below the auxiliary anchor surplus, then a logarithmic cost in
the symbol length when \(s'>s\). Slightly enlarging the groups gives
a further saving on each tail bit. The last sum of corrections then
converges uniformly in the alphabet, giving the required bound on
the cost of one column.
